\section{Proceso metodológico}
El proceso para ambos se compone de tres fases: planificación, conducción y reporte, con diferencias fundamentales en el enfoque de cada fase para cada metodología.
\subsection{Planificación}
Durante la fase de planificación, se define la justificación y objetivos, se formulan las preguntas y se desarrolla un protocolo.\\ 

En el SMS, el protocolo\cite{definicion_sms} define:
\begin{itemize}
    \item Objetivo y preguntas de mapeo (amplias y exploratorias, enfocándose más en cantidad y tipo de estudios que en resultados detallados).
    \item Estrategias de búsqueda (bases de datos, palabras clave, etc.).
    \item Criterios de inclusión y exclusión de fuentes.
    \item Procedimiento de selección.
    \item Esquema de clasificación (define facetas: tema, tipo de contribución, tipo de investigación, contexto).
    \item Extracción de datos y análisis (formularios, plan de análisis de datos).
    \item Reporte (formato de presentación del mapa).\\
\end{itemize}

El protocolo del SLR\cite{kitchenham2009slr} se compone de:
\begin{itemize}
    \item Contexto, objetivo y preguntas de investigación o RQs (preguntas específicas buscando evidencias).
    \item Estrategias de búsqueda (bases de datos, palabras clave, etc.)
    \item Criterios de inclusión y exclusión de fuentes
    \item Procedimiento de selección
    \item Evaluación de calidad (checklists, uso de scores)
    \item Extracción y síntesis de datos (campos de extracción, tipo de síntesis, plan de análisis)
    \item Reporte (cronograma, versionamiento, plan de reporte)\\
\end{itemize}

\subsection{Conducción}
Este es el núcleo de ambas metodologías, aunque ambas tienen una estructura similar general, difieren significativamente en la profundidad y el tipo de datos que se procesan.\\
 
En el SMS, la conducción tiene como características que:
\begin{itemize}
    \item La investigación es mucho más amplia. El objetivo es cubrir toda la producción de un área temática o campo de estudio.
    \item La selección de estudios primarios puede ascender a los miles artículos los cuales se van filtrando a medida que avanza la investigación.
    \item La síntesis de datos es de alto nivel y se enfoca en la categorización y clasificación de los mismos.\\
\end{itemize}

Por otro lado, en la metodología SLR:
\begin{itemize}
    \item La investigación se enfoca en responder las preguntas de la investigación.
    \item La selección de estudios primarios es riguroso para filtrar estudios que no respondan directamente a la hipótesis o pregunta puntual.
    \item La síntesis de datos es profunda, así como cualitativa y cuantitativa.\\ 

\end{itemize}

\subsection{Reporte}
En la fase de reporte se describe el proceso y los resultados obtenidos pero la salida es diferente: una síntesis de evidencia en el SLR, y un mapa de terreno en el SMS.\\

En un SMS, el reporte típicamente utiliza las siguientes secciones para describir el paisaje:
\begin{itemize}
    \item Descripción del proceso.
    \item Presentación del esquema de clasificación.
    \item Resultados de mapeo (tablas y gráficos que muestren frecuencias y distribuciones).
    \item Identificación de tendencias y vacíos.\\
\end{itemize}

Por otro lado, el SLR está orientado a responder cada pregunta de investigación utilizando:
\begin{itemize}
    \item Descripción del proceso.
    \item Características de los estudios incluidos (tablas, scores).
    \item Síntesis de resultados.
    \item Discusión y conclusiones.\\
\end{itemize}